%══════════════════════════════════════════════════════════════════════════════
%─── Section 5: Monte Carlo ──────────────────────────────────────────────────
%══════════════════════════════════════════════════════════════════════════════
\section{Monte Carlo Evidence}
\label{sec:mc}

Theorem~\ref{thm:dependent_conformal_var} is asymptotic and its remainder is
conservative. This section measures what the correction delivers at sample sizes
a practitioner has, under misspecification of a known kind and severity, and
records one thing that does not converge.

\paragraph{Design} Returns follow $r_t = \sigma_t\varepsilon_t$ with a
GARCH$(1,1)$ recursion at $\omega = 10^{-5}$, $\alpha_1 = 0.10$,
$\beta_1 = 0.85$, so Corollary~\ref{cor:garch} applies to every cell. Five
innovation laws are used, in increasing order of tail departure: Normal, which is
correct specification; Student-$t(5)$; Student-$t(3)$; the skewed-$t(3,-0.5)$ of
\citet{hansen1994autoregressive}; and the mixture
$0.95\,\mathcal N(0,1)+0.05\,\mathcal N(0,25)$. All are standardised to unit
variance. The forecaster is given the \emph{true} GARCH parameters and assumes
Normal innovations, so the only quantity varying across the five processes is the
innovation law and the only defect is distributional.

Sample sizes are $T\in\{500,\,1000,\,2000,\,5000,\,10000\}$ with $f_c = 0.70$ and
$\alpha = 0.01$, giving calibration samples from 350 to 7{,}000 observations, and
each cell is replicated 500 times: 12{,}500 replications in total. The shift is
equation~\eqref{eq:qV} throughout.

\subsection{What converges}
\label{sec:mc_converges}

Table~\ref{tab:mc_grid} reports the grid and Figure~\ref{fig:mc_convergence}
plots it.

\begin{table}[htbp]
	\centering
	\caption{Monte Carlo across five data-generating processes and five sample
		sizes, 500 replications per cell, $\alpha = 0.01$, $f_c = 0.70$. All
		processes are GARCH$(1,1)$ with $\omega = 10^{-5}$, $\alpha_1 = 0.10$,
		$\beta_1 = 0.85$; the forecaster uses the true GARCH parameters and
		assumes Normal innovations. Corrected $\hat\pi$ is the post-correction
		violation rate on the test block, averaged across replications; Green is
		the share of replications whose \emph{raw} forecasts fall in the Basel
		green zone. The green zone is $\hat\pi\le 4/250 = 0.016$, and the
		population violation rate of each process is given in the first column
		for comparison against it. The $T = 500$ row of each panel carries a
		small-sample bias in $\qVstat$ of both signs and is reported with that
		qualification (Supplement~S.4.1).}
	\label{tab:mc_grid}
	\footnotesize
	\setlength{\tabcolsep}{4.5pt}
	\begin{tabular}{@{}lc rrrrr rrrrr@{}}
		\toprule
		& population
		& \multicolumn{5}{c}{corrected $\hat\pi$}
		& \multicolumn{5}{c}{raw forecasts in the green zone (\%)} \\
		\cmidrule(lr){3-7}\cmidrule(l){8-12}
		Innovations & $\hat\pi$
		& 500 & 1{,}000 & 2{,}000 & 5{,}000 & 10{,}000
		& 500 & 1{,}000 & 2{,}000 & 5{,}000 & 10{,}000 \\
		\midrule
		Normal (correct) & 0.0100 & 0.0076 & 0.0101 & 0.0100 & 0.0100 & 0.0100 & 83.4 & 80.6 & 92.6 & 98.8 & 100.0 \\
		Student-$t$(5) & 0.0150 & 0.0085 & 0.0100 & 0.0100 & 0.0102 & 0.0098 & 61.6 & 54.2 & 57.8 & 65.0 & 74.2 \\
		Student-$t$(3) & 0.0136 & 0.0084 & 0.0095 & 0.0099 & 0.0098 & 0.0101 & 69.6 & 61.4 & 67.4 & 81.6 & 85.6 \\
		Skewed-$t$(3, $-0.5$) & 0.0233 & 0.0083 & 0.0104 & 0.0105 & 0.0101 & 0.0101 & 36.4 & 18.0 & 10.8 & 3.2 & 0.4 \\
		Mixture of Normals & 0.0125 & 0.0084 & 0.0094 & 0.0098 & 0.0100 & 0.0099 & 68.4 & 68.2 & 79.8 & 90.8 & 96.4 \\
		\midrule
		\multicolumn{12}{@{}l}{\textit{Mean $\qVstat$ (left) and its standard deviation across replications (right)}} \\[2pt]
		Normal (correct) & --- & 0.00165 & 0.00030 & 0.00015 & 0.00008 & 0.00003 & 0.00307 & 0.00192 & 0.00132 & 0.00088 & 0.00064 \\
		Student-$t$(5) & --- & 0.00758 & 0.00459 & 0.00405 & 0.00369 & 0.00372 & 0.00687 & 0.00383 & 0.00268 & 0.00163 & 0.00113 \\
		Student-$t$(3) & --- & 0.00868 & 0.00493 & 0.00424 & 0.00365 & 0.00360 & 0.01093 & 0.00526 & 0.00331 & 0.00200 & 0.00142 \\
		Skewed-$t$(3, $-0.5$) & --- & 0.02112 & 0.01383 & 0.01249 & 0.01220 & 0.01214 & 0.01888 & 0.00843 & 0.00479 & 0.00302 & 0.00222 \\
		Mixture of Normals & --- & 0.01451 & 0.00863 & 0.00727 & 0.00631 & 0.00601 & 0.01274 & 0.00846 & 0.00622 & 0.00388 & 0.00286 \\
		\bottomrule
	\end{tabular}
\end{table}

\begin{figure}[htbp]
	\centering
	\includegraphics[width=\textwidth]{fig_mc_convergence}
	\caption{Monte Carlo across five data-generating processes and five sample
		sizes, 500 replications per cell. Panel (a): the post-correction
		violation rate against the nominal $\alpha = 0.01$. Panel (b): the share
		of replications whose raw forecasts fall in the Basel green zone. Panel
		(c): the standard deviation of $\qVstat$ across replications, on
		logarithmic axes, against a $T^{-1/2}$ reference.}
	\label{fig:mc_convergence}
\end{figure}

Three readings follow from the upper panel of Table~\ref{tab:mc_grid}.

\emph{Coverage is restored across every process.} From $T = 1{,}000$ onward the
corrected violation rate lies within $0.0006$ of nominal for all five processes,
including the skewed-$t$, whose raw rate is $0.0233$ --- more than twice the level
it reports. The residual falls with $T$ and there is no ordering of it by the
severity of the misspecification: a one-parameter shift removes as much of a
severe unconditional error as of a mild one.

\emph{$\qVstat$ converges to the population miscalibration, which does not vanish
with data.} Under correct specification it tends to zero, from $0.00165$ at
$T = 500$ to $0.00003$ at $T = 10{,}000$. Under each misspecification it converges
instead to a non-zero limit, ordered by the severity of the tail departure:
$0.0037$ for Student-$t(5)$, $0.0036$ for Student-$t(3)$, $0.0060$ for the
mixture, $0.0121$ for the skewed-$t$. The quantity being estimated is a property
of the pair of laws, not an artefact of the sample.

\emph{Its variance falls at the parametric rate.} The replication standard
deviation of $\qVstat$ under correct specification falls from $0.00307$ to
$0.00064$ across a twentyfold increase in $T$, a factor of $4.80$ against
$\sqrt{20} = 4.47$. Lemma~\ref{lem:quantile_mixing} carries an additional
$\sqrt{\log n}$; over this range of $T$ the two rates are not separable, and no
attempt is made to separate them.

\paragraph{The small-sample cost of the finite-sample convention} At $T = 500$
the corrected rate is $0.0076$ to $0.0085$ across all five processes, including
the correctly specified one --- systematically conservative rather than nominal.
The mechanism is the conformal index of Section~\ref{sec:theory_lemmas}: at
$n = 350$ the level $\lceil (n+1)(1-\alpha)\rceil/n$ is $0.9943$ rather than
$0.99$, an overshoot of $0.0043$, and $\qVstat$ inherits it. Under correct
specification, where the population shift is zero, the overshoot is the entire
signal: mean $\qVstat$ runs $0.00165,\,0.00030,\,0.00015,\,0.00008,\,0.00003$
across the grid, against overshoots of $0.0043,\,0.0014,\,0.0007,\,0.0003,\,0.0001$.
The correction is upward-biased on a forecaster that needs no correction, by
construction rather than by misspecification, and the bias is $\mathcal O(1/n)$.
Section~\ref{sec:recal} reports the panel consequence.

\subsection{What does not converge}
\label{sec:mc_diverges}

The right-hand panel of Table~\ref{tab:mc_grid} reports the share of replications
whose \emph{raw} forecasts are classified green, and it does not behave like the
left-hand panel.

The Basel green zone is $\hat\pi\le 4/250 = 0.016$
\citep{basel1996}. As $T$ grows the sampling distribution of $\hat\pi$
concentrates on its population value, so the green share runs to $100$\% or to
$0$\% according to which side of $0.016$ that value falls --- a fact verified
replication by replication, with the green share equal to the share of
replications on the correct side of the threshold in all 25 cells. Sample size
therefore sharpens the classification without moving the boundary, and the two
directions are visible in the same table.

The skewed-$t$ process, whose population rate of $0.0233$ clears the boundary, is
driven out: $36.4$\% green at $T = 500$ falls to $0.4$\% at $T = 10{,}000$. The
Student-$t(5)$ process is not. Its population rate is $0.0150$ --- a forecaster
understating its exceedance probability by half --- and it sits inside the
boundary, so the same concentration works the other way: $61.6$\% green at
$T = 500$ rises to $74.2$\% at $T = 10{,}000$. Student-$t(3)$ at $0.0136$ rises
from $69.6$\% to $85.6$\%, and the mixture at $0.0125$ from $68.4$\% to $96.4$\%.

More data makes the traffic light more confident, and for three of the four
misspecified processes here it makes it more confidently green. The classification
converges to a verdict determined by the position of the true miscoverage relative
to a fixed threshold, and not to the miscoverage. What converges to the
miscoverage is $\qVstat$.

This is the identification constraint of Section~\ref{sec:precede} in a setting
where the truth is known by construction, and it is worth separating from a claim
the design does not support. The zone boundary is not miscalibrated: a forecaster
at $\hat\pi = 0.0150$ genuinely produces fewer than four exceedances per
250 days more often than not, which is what the zone is defined to detect. The
point is narrower and is about what a practitioner may infer. A green zone
sustained over a longer record is evidence that the exceedance count is below the
boundary, and it is not evidence that the reported quantile is the quantile it
claims to be --- the two coincide only when the miscoverage is large enough to
cross the boundary. Waiting for more data does not close that gap; it widens the
share of replications in which the gap is invisible.
